\scanpage{6}%
% L1–4; continuation of anti- at the foot of scan 5
podismes $\sigma_i$, on vérifie que pour
\[
I=\{i,j,k\}
\]
on a
\[
\sigma_k=\sigma_i\sigma_j\sigma_i
\]
% L5–8
\deleted{\unsure{d'où|donc}}\unsure{ce|donc} qui \inserted{\unsure{donc|nous|peut}} \unsure{exige|exiger} (en faisant une permutation circulaire, donc
\[
\begin{aligned}
\sigma_i&=\sigma_j\sigma_k\sigma_j
=\sigma_{\deleted{\unsure{i}}j}(\sigma_i\sigma_j\sigma_i)\sigma_j
\qquad\deleted{\text{i.e.}}\\
(\sigma_i\sigma_j)^3&=1
\end{aligned}
\]

% L9–15
On \deleted{constate que $\unsure{G|G'}$ est le sous-groupe présenté} \inserted{voit que $G'$ opérant sur} $I$ donne un iso \deleted{par}
\[
\deleted{\sigma_i^2=1\qquad(i\in I)}
\]
\[
\deleted{\unsure{\sigma}}\qquad G'\xrightarrow{\sim}\unsure{\mathfrak S|\mathfrak G}_I
\]
avec $\sigma_i$ correspondant à la transposition de $I$ d'indice $i$ -- $G'$ est présenté par les gén. $\sigma_i$ \inserted{$(i\in I)$} et les \inserted{$3+6=9$} relations
\[
\sigma_i^2=1\qquad \sigma_k=\sigma_i\sigma_j\sigma_i\qquad
\text{(si $I=\{i,j,k\}$,}
\]

% L16–19
ou encore (pour $i,j\in I$ fixés, $i\ne j$ \deleted{\unsure{$j$|)}}, posant \unsure{ici} $\sigma=\sigma_i,\sigma'=\sigma_j$)
\[
\deleted{\sigma_i^2=\sigma_j^2=}
\]
\[
\sigma^2=\sigma'{}^2=(\sigma\sigma')^3=1\quad.
\]

% L20–21
Quel est le groupe quotient $G/G'$ (d'ordre $6$) ? Notons que \deleted{\unsure{les}}\unsure{la|les} \inserted{donnée} \deleted{\unsure{str.}} de $(S_i)_{i\in\unsure{\tau|T}}$
% archive sheet number: 5

